\chapter{Conclusion}

This chapter reviews the principal findings of the research and evaluates the effectiveness of the proposed \acs{FPGA} Verification Platform. It summarizes the key contributions, examines how the proof-of-concept addresses the initial challenges outlined in the thesis, and discusses the implications for servo-drive \acs{FPGA} validation. Additionally, the chapter proposes potential extensions and improvements to the platform, highlighting its capacity to facilitate future advancements in industrial motor-control verification.

\section{Summary}

This thesis presents the design, development, and validation of a dedicated \acs{FPGA} Verification Platform for servo-drive control boards utilizing the Zynq UltraScale+ \acs{MPSoC} architecture. The growing complexity of \acs{FPGA} modules in industrial motor control, along with limited test coverage in current verification methodologies and the inherent shortcomings of simulation-based approaches, motivated this research. Simulation methods often do not capture real-world timing, analog feedback, or module interactions, thus, in-system validation is essential to ensure reliable performance.

To address these challenges, this thesis introduces a Proof-of-Concept \acs{FPGA}-based verification framework. The system generates realistic input signals, executes automated tests, and evaluates device responses within a closed-loop configuration. Operating independently from the complete motor-control system, it enables black-box testing of individual modules. This approach enhances test coverage, simplifies debugging, and reduces the need for additional hardware.

A comprehensive technical analysis of the \acs{MPSoC} architecture, clocking strategies, verification methodologies, and \acs{AMBA} \acs{AXI} communication protocols formed the foundation for this research. Based on these considerations, three critical \acs{FPGA} modules from the servo-drive control board were selected for verification - Current Sense \acs{SPI} path, Over-Current Protection module, and \acs{PWM} generation module.

For each module, dedicated systems were developed to generate test signals, measure outcomes, and analyze data. The Current Sense verification employed both static and dynamic tests using a DAC80508 analog signal, Processing subsystem for \acs{FFT} analysis, and MATLAB for output validation. The \acs{PWM} verification assessed frequency, duty cycle, and dead time with high precision. The \acs{OVC} verification evaluated event codes and protection behavior across various programmed fault thresholds.

Experimental results demonstrated that all tested modules responded appropriately to synthetic inputs. Measurement errors remained within specified limits, and each test concluded with a "TEST PASSED" outcome. The platform successfully replicated key characteristics of motor-control signals, verified the correct operation of the protection logic, and demonstrated strong agreement between the reference and measured values.

The proposed verification platform addresses the primary challenges outlined in the section \ref{sec:intro:ProblemStatement} (Problem Statement) . Automated test execution and a standalone \acs{MPSoC}-based configuration minimize manual intervention and eliminate the need for repeated reconfiguration, thereby increasing automation. The platform can generate both static and dynamic input signals without physical motors, demonstrating that realistic operating conditions can be fully emulated during \acs{FPGA} verification. By providing each module with an independent test system, the platform enables genuine module-level validation and eliminates the interdependence among \acs{FPGA} modules, such as the need for a \acs{PWM} module during Current Sense Verification. Digital measurement techniques, including \acs{FFT}-based analysis and precise signal verification, replace manual estimation, thereby improving accuracy and repeatability. Experimental results indicate that this approach overcomes the principal limitations of previous methods and provides a more scalable, flexible, and reliable solution for \acs{FPGA} module verification.

In summary, the verification platform demonstrates that a low-power, fully programmable, \acs{MPSoC}-based \acs{FPGA} verification approach can effectively support both analog and digital testing. This Proof-of-Concept establishes a robust foundation for the development of future scalable verification systems in industrial servo-drive applications.

\section{Future Work}

As servo-drive control systems become more complex, the proposed \acs{FPGA} verification platform can be further enhanced and expanded. Future developments should prioritize improved communication, broader verification coverage, and enhanced usability through increased automation.

A significant upgrade for the platform involves replacing the current \acs{UART} communication with an Ethernet-based interface. This modification would enable faster data transfer, reduced latency, and more reliable connections between the PC and the Verification \acs{MPSoC}, particularly when handling large datasets or real-time measurements.

The platform may also be extended to verify additional \acs{FPGA} modules on the control board. Future enhancements should include encoder interfaces, brake-control logic, analog measurement paths, Ethernet \acs{TSN} features, and digital I/O modules. Incorporating these elements would facilitate more comprehensive board-level verification.

Subsequent versions of the platform may incorporate \acs{FPGA}-accelerated \acs{HIL} models to simulate complete motor-control systems or power electronics. This approach would enable realistic closed-loop testing without the need for physical motors or power hardware, aligning the proof of concept with contemporary \acs{HIL} and \acs{FIL} verification methodologies.

Usability may be enhanced by developing a high-level software interface for automated testing, data logging, and result visualization. The addition of dashboards and structured reports would streamline verification processes and improve accessibility for larger engineering teams.

Integrating the verification platform into continuous integration and continuous deployment (CI/CD) pipelines would enable automated regression testing and hardware validation throughout \acs{FPGA} development. This integration would reduce development risks and support sustained design reliability.

In conclusion, these enhancements have the potential to transform the current proof-of-concept into a robust, scalable, automated, and industry-ready \acs{FPGA} verification platform suitable for modern servo-drive control systems.

%The presented verification platform presents multiple opportunities for enhancement and expansion. As servo-drive control systems become more complex, subsequent updates can improve the platform’s performance, scalability, and usability. The following outlines potential directions for future research and development.
%
%A significant upgrade for the platform involves replacing the current UART communication with an Ethernet-based interface. This modification would enable faster data transfer, reduced latency, and more reliable connections between the PC and the Verification MPSoC, particularly when handling large datasets or real-time measurements.
%
%The platform's verification scope can be expanded to include additional FPGA modules on the control board. Incorporating support for encoder evaluation, brake-control logic, analog measurement paths, Ethernet-TSN features, and digital I/O modules would significantly enhance test coverage and facilitate comprehensive board-level verification.
%
%Platform improvement can also be achieved by transitioning from the current analog stimulus path to a fully digital approach. Employing digitally generated test signals or bitstream based DAC emulation would reduce external dependencies, increase test repeatability, and simplify hardware setup, while maintaining precise control over test signals.
%
%Future iterations of the platform may incorporate FPGA-accelerated hardware-in-the-loop models capable of simulating complete motor-control systems or power-electronic stages. This enhancement would enable realistic closed-loop testing without the need for physical motors or power hardware, aligning the proof of concept with contemporary HIL and FIL verification methodologies.
%
%To improve usability, a high-level software interface could be developed to support automated test execution, logging, and visualization. The addition of dashboards and comprehensive reporting tools would streamline verification tasks and enhance accessibility for larger engineering teams.
%
%Integrating the platform into continuous integration and continuous deployment (CI/CD) pipelines would facilitate ongoing verification of FPGA modules. Automated regression testing, nightly builds, and hardware-based functional testing would mitigate development risks and maintain FPGA design reliability throughout iterative changes.
%
%In conclusion, the proposed improvements have the potential to transform the current proof-of-concept into a robust, scalable, and fully automated verification environment. These enhancements would facilitate support for additional FPGA modules, expedite development cycles, improve design reliability, and streamline maintenance within modern servo-drive control systems.
